\chapter{AKSZ Formalism and the Classical BV Construction}

\section{Introduction and Motivation}

This chapter introduces the Alexandrov--Kontsevich--Schwarz--Zaboronsky
(AKSZ) construction for classical field theories with
a $(k-1)$-shifted symplectic target, following
\cite{AKSZ97,CF01,CF07,PTVV13}.
The formalism generalizes the classical Batalin–Vilkovisky (BV) method
and provides a canonical route to constructing solutions of the
classical master equation. The approach is intrinsically geometric,
relying on mapping stacks, graded manifolds,
and homological vector fields of degree $+1$.

From the perspective of RSVP, the AKSZ framework is essential.
RSVP fields will be realized as maps
\[
\mathcal{F} \colon \mathsf{T}[1]M \longrightarrow \mathcal{X},
\]
where $\mathcal{X}$ is the derived moduli stack of
lamphrodynamic triples $(\Phi,\vec v,S)$ constructed in Chapter~6,
equipped with a canonical $(-1)$-shifted symplectic structure.
The AKSZ recipe then produces a BV action functional $S_{\text{BV}}$
satisfying $\{S_{\text{BV}},S_{\text{BV}}\}=0$, thereby yielding a
consistent classical theory in the BV sense.

This chapter develops the required formal background, establishes the
general AKSZ construction, and explains the classical BV complex.
The final section illustrates the formalism with the
example of three-dimensional Chern--Simons theory.

\medskip

\noindent{\bf Pedagogical note.}
The first half of this chapter is written from the perspective of a
physicist learning derived geometry; the second half assumes familiarity
with shifted symplectic structures, homological vector fields,
and the BV spectral sequence.


\section{Graded Manifolds and Homological Vector Fields}

Let $\mathcal{M}$ be a $\mathbb{Z}$-graded supermanifold (or more
generally, a derived stack locally modeled by such).
A \emph{homological vector field}
is a degree $+1$ vector field $Q$ on $\mathcal{M}$ satisfying
\[
Q^2 = \tfrac12[Q,Q]=0.
\]
The pair $(\mathcal{M},Q)$ is called a \emph{$Q$-manifold}.
The condition $Q^2=0$ implies that $Q$ determines a differential on
the algebra of functions $\mathcal{O}_{\mathcal{M}}$, and hence gives
rise to a cochain complex
$(\mathcal{O}_{\mathcal{M}},Q)$.

\begin{definition}
A \emph{$Q$-mapping stack} from $\mathsf{T}[1]M$ to $\mathcal{X}$
is a mapping stack
\[
\mathrm{Map}(\mathsf{T}[1]M,\mathcal{X})
\]
equipped with a homological vector field induced from the
vector fields on both source and target, in particular
the de Rham differential on $\mathsf{T}[1]M$.
\end{definition}

\begin{remark}
Physically, the degree shift encodes ghost number.
Coordinates of degree $0$ describe classical fields,
degree $-1$ antifields, degree $+1$ ghosts.
\end{remark}


\section{Shifted Symplectic Structures and the BV Bracket}

Let $(\mathcal{X},\omega)$ be a $k$-shifted symplectic derived stack.
Thus $\omega$ is a closed, non-degenerate $2$-form of cohomological
degree $k$.
When $k=-1$, the induced Poisson bracket has cohomological degree $+1$,
and is known as the BV bracket.

\begin{definition}
Given a $(-1)$-shifted symplectic structure $\omega$,
the \emph{BV bracket} $\{-,-\}$ on
$\mathcal{O}_{\mathcal{X}}$ is defined by
\[
\{f,g\} := \iota_{X_f}\iota_{X_g}\omega,
\]
where $X_f$ is the Hamiltonian vector field determined by $f$.
\end{definition}

\begin{proposition}
The BV bracket has the following properties:
\begin{enumerate}
\item $\{f,g\} = -(-1)^{(|f|+1)(|g|+1)}\{g,f\}$,
\item $\{f,gh\} = \{f,g\}h + (-1)^{(|f|+1)|g|}g\{f,h\}$,
\item $(-1)^{(|f|+1)(|h|+1)}\{f,\{g,h\}\} + \text{cyclic perms} = 0$.
\end{enumerate}
\end{proposition}

\begin{proof}[Proof (sketch)]
The identities follow from the graded symplectic structure and the
Cartan formula. A complete proof may be found in
\cite{Schwarz93,CattaneoFelder02,PTVV13}.
\end{proof}


\section{AKSZ Data and the Classical Master Equation}

Let $M$ be a smooth oriented manifold of dimension $d$,
and let $(\mathcal{X},\omega)$ be a $(d-1)$-shifted symplectic target
with a homological vector field $Q$ preserving $\omega$.
The AKSZ construction associates to this data a BV action on the
mapping stack
\[
\mathcal{F} := \mathrm{Map}(\mathsf{T}[1]M,\mathcal{X}),
\]
with an induced $(-1)$-shifted symplectic structure.
Concretely, the AKSZ action is given by
\begin{equation}
S_{\mathrm{AKSZ}} = \int_{\mathsf{T}[1]M}
\left( \langle \Phi^{\ast}\theta, d_M \Phi\rangle
    + \Phi^{\ast}\alpha \right),
\label{eq:AKSZ-action}
\end{equation}
where $\theta$ is a potential of $\omega$ and $\alpha$ is a Hamiltonian
function satisfying $Q=\{\alpha,-\}$.

\begin{theorem}[AKSZ classical master equation]
The functional $S_{\mathrm{AKSZ}}$ satisfies
\[
\{S_{\mathrm{AKSZ}},S_{\mathrm{AKSZ}}\}=0.
\]
\end{theorem}

\begin{proof}[Proof (sketch)]
The construction of $S_{\mathrm{AKSZ}}$ is arranged so that
the BV bracket encodes the graded symplectic form on the mapping space,
and the homological vector field encodes the differential on
$\mathsf{T}[1]M$.
The fact that $Q$ preserves $\omega$ and squares to zero implies
that $S_{\mathrm{AKSZ}}$ is a solution of the classical master equation.
For complete details see \cite{AKSZ97,CF01,CF07,PTVV13}.
\end{proof}

\section{Mapping Stacks and the Induced Shifted Symplectic Form}

Let $(\mathcal{X},\omega)$ be a $(d-1)$-shifted symplectic derived stack
in the sense of \cite{PTVV13}, and consider a compact oriented
$d$-manifold $M$. The AKSZ construction uses the internal hom
$\mathrm{Map}(\mathsf{T}[1]M,\mathcal{X})$, viewed as a derived mapping
stack in the $\infty$-topos of derived stacks. The tangent complex
is obtained by the standard formula
\[
T_{\Phi}\mathrm{Map}(\mathsf{T}[1]M,\mathcal{X})
\simeq \mathbf{R}\Gamma(\mathsf{T}[1]M, \Phi^{\ast}T_{\mathcal{X}}),
\]
which identifies tangent vectors at a field configuration $\Phi$
with sections of the pullback of the tangent complex of $\mathcal{X}$.

\begin{theorem}[PTVV]
\label{thm:PTVV-induced}
Let $(\mathcal{X},\omega)$ be $(d-1)$-shifted symplectic
and let $M$ be an oriented $d$-manifold.
Then the mapping stack $\mathrm{Map}(\mathsf{T}[1]M,\mathcal{X})$
carries a canonical $(-1)$-shifted symplectic structure
$\omega_{\mathrm{AKSZ}}$ induced by integration along $M$.
\end{theorem}

\begin{proof}[Proof (sketch)]
The degree shift arises from the Thom form of $M$ and the graded geometry
of $\mathsf{T}[1]M$. Integration reduces degree by $d$, and combined with
the $(d-1)$-shift of $\omega$, yields $-1$. The construction and closedness
are established in \cite{PTVV13}.
\end{proof}

\begin{remark}
This result is fundamental:
any $(d-1)$-shifted symplectic target yields a $(-1)$-shifted symplectic
BV theory on the mapping stack. In particular, the RSVP target constructed
in Chapter~6 will have $d=3$, hence a $(-1)$-shifted BV theory in
four spacetime dimensions.
\end{remark}


\section{BV Fields, Ghosts, and Antifields}

Given the AKSZ mapping stack
$\mathcal{F}=\mathrm{Map}(\mathsf{T}[1]M,\mathcal{X})$,
the derived geometry automatically organizes field components by ghost
number. More precisely, a degree-$k$ coordinate on $\mathcal{X}$ pulls
back to a degree-$(k+i)$ coordinate on $\mathcal{F}$, where $i$ labels
the internal grading of $\mathsf{T}[1]M$.
Classical fields occur at total degree $0$, ghosts at $+1$, antifields at
$-1$, and higher ghosts at degrees $>+1$.

\begin{definition}
The \emph{BV complex} associated to $\mathcal{F}$ is the cochain complex
$(\mathcal{O}_{\mathcal{F}},Q_{\mathrm{AKSZ}})$, where
$Q_{\mathrm{AKSZ}}$ is the homological vector field induced from both the
target differential and the de~Rham differential on the source.
\end{definition}

\begin{remark}
In classical physics language, we may write fields schematically as
\[
\{ \text{fields } \phi,\;\text{ghosts } c,\;\text{antifields }
\phi^{\ast},\;\text{anti-ghosts } c^{\ast},\ldots \},
\]
with degrees organized by the AKSZ grading.
\end{remark}


\section{The Classical BV Action and the Master Equation}

Let $(\mathcal{X},\omega)$ be $(d-1)$-shifted symplectic, and let
$\theta$ be a potential of $\omega$ (i.e.\ $\omega = d\theta$).
Assume $\mathcal{X}$ carries a Hamiltonian function $\alpha$ of degree
$d$ such that the associated Hamiltonian vector field $Q=\{\alpha,-\}$
coincides with the homological vector field on $\mathcal{X}$. Then
the AKSZ action \eqref{eq:AKSZ-action} satisfies the BV classical
master equation
\[
\{S_{\mathrm{AKSZ}},S_{\mathrm{AKSZ}}\} = 0.
\]
This is the precise sense in which AKSZ produces a BV theory.

\begin{proposition}
The BV bracket associated to $\omega_{\mathrm{AKSZ}}$ coincides with
the variational Schouten bracket on functionals over the mapping stack.
\end{proposition}

\begin{proof}[Proof (sketch)]
This is essentially the compatibility of the PTVV bracket with the internal
Hom structure and the universal property of the mapping stack. See
\cite{PTVV13,CF07}.
\end{proof}


\section{Example: Three-Dimensional Chern--Simons Theory}

We illustrate the general AKSZ recipe by recovering three-dimensional
Chern--Simons gauge theory, following \cite{AKSZ97,CF01,CF07}.
Let $M$ be a $3$-dimensional oriented manifold and let
$\mathfrak{g}$ be a finite-dimensional Lie algebra with invariant bilinear
form $\langle\cdot,\cdot\rangle$. Consider the graded vector space
\[
\mathfrak{g}[1],
\]
concentrated in degree $1$, i.e.\ coordinates have ghost number $1$.
Equip $\mathfrak{g}[1]$ with the homological vector field
\[
Q(x) = \tfrac12 [x,x],
\]
which is of degree $+1$ and squares to zero by the Jacobi identity.

\begin{proposition}
The space $\mathfrak{g}[1]$ carries a canonical
$2$-shifted symplectic form (hence $(d-1)=2$), given by the bilinear form
of degree~$2$ induced from $\langle\cdot,\cdot\rangle$.
\end{proposition}

\begin{proof}[Proof (sketch)]
The invariant bilinear form provides a non-degenerate pairing of appropriate
degree, and closedness follows from the Lie algebra structure. See
\cite{AKSZ97}.
\end{proof}

Applying the general AKSZ recipe with $d=3$ and target $\mathfrak{g}[1]$,
the mapping stack $\mathrm{Map}(\mathsf{T}[1]M,\mathfrak{g}[1])$ obtains
a $(-1)$-shifted symplectic structure.
Let $A$ denote the superfield obtained from the pullback of the graded
coordinate on $\mathfrak{g}[1]$; its degree-$0$ component is an ordinary
connection $1$-form on $M$, and higher-degree components encode ghosts
and antifields.

The AKSZ action takes the familiar form
\begin{equation}
S_{\mathrm{CS}} = \int_M \left\langle A \wedge dA
+ \tfrac13 A \wedge [A,A] \right\rangle,
\end{equation}
where products and brackets are understood in the graded sense.
One verifies that $\{S_{\mathrm{CS}},S_{\mathrm{CS}}\}=0$ using the
graded Jacobi identity and invariance of $\langle\cdot,\cdot\rangle$.

\begin{remark}
The usual Chern--Simons action is the restriction of $S_{\mathrm{CS}}$
to the classical degree-zero part of the AKSZ field.
All ghosts, antifields, and higher-degree components arise automatically
from the AKSZ mapping stack, rather than being introduced by hand.
\end{remark}

\medskip

\noindent\textbf{Forward link to RSVP.}
In Chapter~7 we apply the AKSZ recipe to the
$3$-shifted symplectic target $\mathrm{Plenum}$ of the lamphrodynamic
triples $(\Phi,\vec v,S)$, constructed explicitly in Chapter~6.
Because $\mathrm{Plenum}$ is $(3-1)=2$-shifted symplectic, the associated
mapping stack over $\mathsf{T}[1]M$ in four dimensions acquires a canonical
$(-1)$-shifted BV symplectic structure, producing the classical
RSVP BV action.


